\section{Algebraic tails and full incoming contact}\label{sec:tails}

In this section we prove that the two ends of the minimum have fixed
apartments and the required algebraic labels. We work over $L$ when
choosing bases for norms whose splitness is not yet known. The scalar
extension result identifies every conclusion about the resulting norms
with a conclusion about the original minimum over $K$.

\subsection{The outgoing tail}

\begin{proposition}\label{prop:outgoingtail}
For all sufficiently large $u$, each vertex norm of $c_*$ lies in one
fixed $K$-apartment. In its fixed basis its logarithmic weights satisfy
\begin{equation}\label{eq:outgoingequation}
 R(u)\tan w_{vi}'(u)=\epsilon_v a\sinh u+p_{vi},
 \qquad \epsilon_A=-1,\quad\epsilon_C=1.
\end{equation}
The identity arrow is strictly contracting. The associated positions are
$x_{vi}(t)=\epsilon_v at+p_{vi}$, and their labels are pure sums of $A$
or $C$. The full normalized norm at each vertex converges to its
prescribed limit at rate $O(e^{-u})$.
\end{proposition}

\begin{proof}
Fix $J_0>J_*$. The tail estimate in the proof of
Theorem~\ref{thm:relaxedminimum} bounds the distance of every member
$p$ of a minimizing sequence with $J(p)\le J_0$ from the outgoing
reference norms by
\[
 d_\infty(p_v(u),r_v(u))\le\eta(u),\qquad
 \eta(u)=2\left(J_0\int_u^\infty\sech s\,\dd s\right)^{1/2}.
\]
The same holds for $c_*$. Choose $M$ with $h_A^+\le e^M h_C^+$,
and choose $U$ so large that $2F(U)>M+2\eta(U)+1$.
Any pair in these two tubes, at any $u\ge U$, satisfies
\begin{equation}\label{eq:tailordertube}
 h_A\le e^{M+2\eta(u)-2F(u)}h_C\le e^{-1}h_C.
\end{equation}

Over $L$, choose separately at each vertex a common basis for
$(c_*)_v(U)$ and $h_v^+$. Let $\mathcal A_v$ be its apartment and
$P_v$ the projection of Lemma~\ref{lem:blockprojection}. It fixes
both these norms and the entire reference curve $r_v$.
Choose a global polygonal minimizing sequence $p_m$, with cutoffs
$U_m\to\infty$, whose velocities have a common strict margin on
$[U,U+1]$. Project its nodes on $[U,U_m]$ by $P_v$, and interpolate
on the same partition. Each projected node stays in its reference
tube, so \eqref{eq:tailordertube} makes the projected pairs ordered.
Their interpolants are ordered as well. Although the two projections
can differ, the tube estimate supplies the order required here.
They preserve both determinant logarithms, hence all endpoint
corrections in the renormalized action.

Let $D_m\ge0$ be the decrease in $E_{[U,U_m]}$ under this node
projection. The original norm at $U$ need not be fixed by $P_v$.
Its distance $\delta_m$ from the respective apartments tends to
zero. Choose $\tau_m\downarrow0$ with $\delta_m/\tau_m\to0$,
refine at $U+\tau_m$, and bridge from the original pair $p_m(U)$
to the projected pair at $U+\tau_m$. The bridge has velocities
bounded by $b+\delta_m/\tau_m$ for a fixed $b<\pi/2$.
It is ordered, fixes the original initial endpoint, and costs
$O(\tau_m)$ more than the projected short interval. Leave the
curve unchanged before $U$. This gives a member of $\calQ_L$ with
\[
 J_*\le J(\widetilde p_m)\le J(p_m)-D_m+O(\tau_m).
\]
Thus $D_m\to0$.

Write $a_{vj}=d_2(p_{m,v}(u_j),\mathcal A_v)$. The departure
estimate of Lemma~\ref{lem:projectionloss} gives
\[
 D_m\ge\frac18\sum_{v,j}\frac{W_j}{\Delta_j^2}
                       (a_{v,j+1}-a_{vj})^2.
\]
The departure at the final cutoff is zero. Weighted
Cauchy--Schwarz, summed backwards from that endpoint, therefore gives
\[
 \sup_{u\ge U,v}d_2(p_{m,v}(u),\mathcal A_v)^2
 \le8\pi D_m,
\]
since $\sum\Delta_j^2/W_j\le\int_U^\infty\sech u\,\dd u\le\pi$.
Convexity of distance to an apartment extends the node bound to the
segments. Closedness and uniform convergence imply that $(c_*)_v$
lies in $\mathcal A_v$ on the whole tail.

Its apartment weights are locally $C^{1,1}$. Indeed, in a fixed
apartment the quadratic chord estimate is the Euclidean estimate
for the unsorted weight coordinates. The dyadic argument of
Proposition~\ref{prop:spectralvelocities}, applied to their coordinate
increments, gives equal forward and backward derivatives and a
locally Lipschitz derivative. The strict local velocity margin
persists. Thus both the weights and their smooth compactly
supported variations are covered by Lemma~\ref{lem:regularaction}. The strict order
\eqref{eq:tailordertube} allows either sign of a small compactly
supported variation of each individual weight. Local minimality
gives, distributionally,
\[
 (R\tan w_{vi}')'=\epsilon_v aR.
\]
Integration proves \eqref{eq:outgoingequation}. Its right side makes
the weights smooth. If $b_{vi}$ is the logarithmic weight of $h_v^+$
in this basis, the full end limit and the elementary bound
$|\arctan(x+y)-\arctan x|\le |y|$ give
\begin{equation}\label{eq:outgoingrate}
 |w_{vi}(u)-\epsilon_vF(u)-b_{vi}|
 \le |p_{vi}|\int_u^\infty\sech s\,\dd s
 \le2|p_{vi}|e^{-u}.
\end{equation}
The reduction of the strictly contracting arrow is zero. The
straight branches consequently have pure $A$ or $C$ labels,
subdivided at their finitely many intersections.

It remains to descend the fixed apartments. The next lemma applies
to $e^{-\epsilon_vF(u)}(c_*)_v(u)$, whose limit is the split
$K$-norm $h_v^+$. It gives a fixed $K$-basis on a shorter tail.
Canonical extension then identifies its restriction with the
original $K$-minimum. Weights differ only by fixed offsets, so their
equations, labels and convergence rates are unchanged.
\end{proof}

\begin{lemma}\label{lem:taildescent}
Let $h$ be a split $K$-norm. Suppose an $L$-norm curve $k(s)$ lies
in one fixed apartment and converges to the extension of $h$ as
$s\to\infty$. On a sufficiently late tail, $k$ is diagonal in one
fixed $K$-basis.
\end{lemma}

\begin{proof}
The limiting norm is diagonal in the same $L$-basis. Rescale its
vectors by real powers of $T$ to make that basis orthonormal for
$h$. Its weights $z_i(s)$ for $k(s)$ then tend to zero. Choose an
orthonormal $K$-basis $f$ for $h$. The change-of-basis matrix
$P$ lies in $\mathrm{GL}_n(\mathcal O_L)$. Its residue is an
invertible complex matrix; let $P_0$ be its constant lift. Write
$P_0^{-1}P=I+Q$. There are only finitely many entries, so all nonzero
entries of $Q$ have valuation at least some $\delta>0$.
The inverse has the same property, by the adjugate formula modulo
the ideal of elements of valuation at least $\delta$.

For large $s$, the weight spread $\max z_i(s)-\min z_i(s)$ is
less than $\delta$. Both $I+Q$ and its inverse are then
nonexpansive for the weighted maximum norm: every off-diagonal
weighted entry has negative logarithm, and the diagonal entries
are units. Thus replacing $P$ by $P_0$ preserves that norm exactly.
The basis $fP_0$ is the required fixed $K$-basis. This uses the
fixed original apartment and the split limit; convergence alone
would not give the conclusion.
\end{proof}

\subsection{The equation on a free interval}
The same scalar equation holds on every interval where the norm
order is uniformly strict, without any assumption about its
location relative to the ends.

\begin{proposition}\label{prop:freeinterval}
Suppose $I=[l,r]$ satisfies $(c_*)_A\le e^{-\kappa}(c_*)_C$
for some $\kappa>0$. Over $L$, choose at each vertex any common
basis for the two endpoint norms. The minimum on $I$ lies in
these fixed apartments and satisfies \eqref{eq:outgoingequation}
for suitable constants $p_{vi}$.
\end{proposition}

\begin{proof}
Let $a_{vi},b_{vi}$ be the endpoint weights. There is a unique
$p_{vi}\in\RR$ such that
\[
 \int_l^r\arctan(\epsilon_v a\tanh u+p_{vi}\sech u)\,\dd u
       =b_{vi}-a_{vi}.
\]
Indeed the left side is strictly increasing with limits
$\pm\pi|I|/2$, and the strict local margin puts the right side
between them. Integrate these velocities from $a_{vi}$ to obtain
a smooth candidate $z$; initially it need not be ordered.

Project any ordered comparison polygon separately onto the two
endpoint apartments. Its kinetic cost decreases and its determinant
potential is unchanged. In those apartments, the supporting
inequality for $\ell$, integrated using
$(R\tan w_{vi}')'=\epsilon_v aR$, proves that its action is at
least the differential action $E(z)$. The fixed basis-volume
constants in the determinant logarithms cancel. This comparison
does not require the projected pair or $z$ to be ordered.
It proves $E_I^*\ge E(z)$.

Take exact-endpoint ordered polygonal almost-minima $p_m\to c_*|I$
and sample $z$ on their refining partitions, obtaining $z_m$.
For large $m$, $p_{m,A}\le e^{-\kappa/2}p_{m,C}$, whereas
$z_{m,A}\le e^M z_{m,C}$ for some finite common $M\ge0$.
Choose $0<\theta<1$ with
$-(1-\theta)\kappa/2+\theta M\le0$. Interpolate the two
polygons' nodes by $\theta$, then join those nodes geodesically.
The resulting polygon $r_m$ is ordered, has the original endpoints,
and satisfies
\[
 E_I^*\le E(r_m)\le(1-\theta)E(p_m)+\theta E(z_m).
\]
Taking limits proves $E(z)=E_I^*$ and $E(r_m)\to E_I^*$.
Local stability makes $r_m\to c_*|I$. Its nodewise limit is
$[c_*(u),z(u)]_\theta$, where brackets denote canonical geodesic
interpolation. Refining meshes identify the same limit between
nodes. Thus $c_*=[c_*,z]_\theta$, which forces $c_*=z$.
\end{proof}

\subsection{Small incoming speeds}
We need equality of the two incoming norms at a finite time.
Contact of only the smallest relative gap would not suffice.

\begin{lemma}\label{lem:smallincoming}
For $s\to-\infty$, the minimum and a global minimizing sequence
can be chosen to have spectral metric speeds bounded by
$O(\cosh(s)^{-1/4})$ on $I_s=[s-1,s+1]$.
For the sequence this is a bound on every segment velocity.
\end{lemma}

\begin{proof}
Refine the time-allocation estimate using a positive weight $\rho$
on an interval of length $H$, with
$0<m\le\rho\le M$ and $|\rho'|\le S$. If the action is at most
$E_0$, the optimized allocation satisfies
\[
 \sum_j\Delta_jL_j\le\frac{2E_0}{m},\qquad
 \sum_j\Delta_j\sum_i v_{ji}^2\le\frac{4E_0}{m}.
\]
Consequently
\[
 \sum_j\Delta_j|\kappa_j|\le\sqrt{8nHE_0/m}.
\]
In \eqref{eq:timeallocationidentity}, the two potential endpoint
weights of each segment sum to its integrated weight. Summing the
kinetic and potential errors therefore gives the improved estimate
\[
 \sum_j|Q_j-Q_{j-1}|
 \le\frac{2SE_0}{m}+qM\sqrt{8nHE_0/m}.
\]
Some segment has $L_j\le B=2E_0/(mH)$. Put
$b_0=\arccos(e^{-B})$ and
\[
 A(E_0)=Mn b_0\tan b_0+\frac{2SE_0}{m}
                         +qM\sqrt{8nHE_0/m}.
\]
Every segment then has $Q_j\le A(E_0)$. Since
$f(v)=v\tan v-\ell(v)\ge v^2/2$,
\[
 |v_{ji}|\le2\sqrt{A(E_0)/m}=O(E_0^{1/4}).
\]
The inverse of $f$ also retains a strict $\pi/2$ margin when
$E_0$ is not small. None of these bounds depends on the number of
nodes.

Fix a global minimizing sequence with $J\le J_0$. On $I_s$,
$s\le-2$, normalize the weight to $\rho_s(u)=R(u)/R(s)$.
Then $e^{-1}\le\rho_s\le e$, $|\rho_s'|\le e$, and
$H=2$. Nonnegativity of the negative-half integrand bounds its
normalized action by $E_0=J_0/R(s)$. Reallocate time on $I_s$,
keeping the endpoints and the exterior fixed. The global action
decreases, so the sequence still converges uniformly to $c_*$.
The preceding bound applies to the sequence and passes to the
limit.
\end{proof}

\begin{proposition}\label{prop:fullcontact}
There is a finite time $U$ at which $(c_*)_A(U)=(c_*)_C(U)$.
\end{proposition}

\begin{proof}
Take $I_s=[s-1,s+1]$ on the negative half-line, and let
$\eta(u)=1-|u-s|$ there, extended by zero outside. Choose a
polygonal minimizing sequence whose velocities on $I_s$ are
bounded by $b_s=O(R(s)^{-1/4})$. For any such polygon $p$, apply
the time-dependent operation
\[
 \widetilde p(u)=S_{\varepsilon\eta(u)}(p_A(u),p_C(u)),
\]
where $S$ is the balanced operation in \eqref{eq:shrinkandcap}.
This variation has the original full endpoints and reference tails.
Lemma~\ref{lem:finitecommonbases} gives a finite common-basis
subdivision of every original segment. Refine at $s$ and at all
crossings of the affine gaps $g_i=c_i-a_i$ with
$\varepsilon\eta$. On each remaining piece the new weights are
\[
 (a_i+\varepsilon\eta/2,c_i-\varepsilon\eta/2)
       \quad\hbox{if }g_i>\varepsilon\eta,
 \qquad
 ((a_i+c_i)/2,(a_i+c_i)/2)
       \quad\hbox{if }g_i<\varepsilon\eta.
\]
They agree on equality intervals. Thus the new curve is a finite
piecewise-apartment member of $\calQ$, remains ordered, and has
velocities at most $b_s+\varepsilon/2<\pi/2$ for small
$\varepsilon$.

Write $b=b_s$ and choose $b<b'<\pi/2$. Taylor's inequality,
with $L_\ell=\sec^2 b'$, bounds the kinetic increase for each
unclipped pair by
\[
 \frac{\varepsilon}{2}\tan b\,|\eta'|
       +\frac{L_\ell\varepsilon^2}{8}|\eta'|^2.
\]
On a clipped pair Jensen's inequality gives no increase. The
determinant gap drops by $\sum_i\min(g_i,\varepsilon\eta)$.
If the largest gap is at least $\gamma>\varepsilon$ throughout
$I_s$, this drop is at least $\varepsilon\eta$. Hence
\begin{equation}\label{eq:tentdecrease}
 J(\widetilde p)-J(p)
 \le-\varepsilon D_s(b)
       +\frac{nL_\ell\varepsilon^2}{8}\int_{I_s}R|\eta'|^2,
 \qquad
 D_s(b)=q\int_{I_s}R\eta-\frac n2\tan b\int_{I_s}R|\eta'|.
\end{equation}
The elementary integrals are
\[
 \int_{I_s}R\eta=2(\cosh1-1)\cosh s,\qquad
 \int_{I_s}R|\eta'|=\int_{I_s}R|\eta'|^2=2\sinh1\cosh s.
\]
It follows that $D_s(b)>0$ precisely when
$n\tan b<\sqrt3\tanh(1/2)$. Choose $s$ sufficiently negative
for this inequality to hold.

If the two norms of $c_*$ were different at every time in $I_s$,
their continuous largest gap would have a positive minimum
$\gamma_*$. Uniform convergence gives a largest gap at least
$\gamma_*/2$ for all sufficiently late polygons. Fix
$\varepsilon>0$, independent of the polygon, below this bound
and small enough that the quadratic term in \eqref{eq:tentdecrease}
is less than $\varepsilon D_s(b)/2$. The transformed sequence
then lowers $J$ by a fixed positive amount, contradicting its
convergence to $J_*$. Thus full equality occurs at a finite time.
\end{proof}

\subsection{The incoming tail}

\begin{proposition}\label{prop:incomingtail}
The full-contact set is $(-\infty,U_*]$ for some finite $U_*$.
On a sufficiently negative tail the two norms coincide in one
fixed $K$-apartment, with weights satisfying
\begin{equation}\label{eq:incomingequation}
 R(u)\tan w_i'(u)=p_i.
\end{equation}
Their actual common norm converges to $h_B$ at rate $O(e^u)$.
The positions are horizontal lines $x_i(t)=p_i$, with pure
semistable $B$ labels.
\end{proposition}

\begin{proof}
Suppose full contact holds at $U$. Project a pair to its common
midpoint, $P_D(h_A,h_C)=(m(h_A,h_C),m(h_A,h_C))$.
Joint convexity of spectral costs shows that this decreases the
total kinetic action of node polygons. It is ordered and removes
the entire nonnegative determinant potential.

Choose a global polygonal minimizing sequence with a common
strict velocity margin on $[U-1,U]$. Apply $P_D$ to all nodes
from its negative cutoff to $U$. Its action decrease $D_m$
satisfies
\[
 D_m\ge q\int_{-U_m}^U R(u)g(p_m(u))\,\dd u\ge0.
\]
The displacement of the final pair tends to zero. As in the
outgoing proof, repair the last interval of length $\tau_m$,
with endpoint displacement $o(\tau_m)$, to the original pair
$p_m(U)$. The bridge remains ordered and has a strict margin.
It changes action by $O(\tau_m)$. Keeping the exterior fixed
and using the truncated renormalization identity gives
$J_*\le J(p_m)-D_m+O(\tau_m)$, even when $U>0$.
Thus $D_m\to0$. Uniform convergence on every compact interval
forces $g(c_*)=0$ at all times before $U$. Ordered norms with
zero determinant gap are equal, since all their relative gaps
are nonnegative. Contact therefore propagates to the entire past.
Nonemptiness follows from Proposition~\ref{prop:fullcontact};
closedness and strict outgoing order identify its stated form.

Over $L$, choose a common basis for $h_B$ and the common norm
at $U_*$. Apply its apartment projection at both vertices on
the negative tail of a minimizing sequence. It preserves order
and determinant gaps. Repair the last endpoint as above.
The projection loss tends to zero. Lemma~\ref{lem:projectionloss},
now summed forwards from zero departure at the negative reference
endpoint, bounds the squared supremum departure by $8\pi$ times
that loss. The limiting common norm consequently lies in this
one apartment on $(-\infty,U_*]$.

As in the outgoing proof, the quadratic chord estimate makes
the fixed-apartment weights locally $C^{1,1}$. Vary each weight
equally at both vertices. The pair
stays equal, its potential is zero, and both signs of the variation
are admissible. Local minimality and
Lemma~\ref{lem:regularaction} give
$(R\tan w_i')'=0$, proving \eqref{eq:incomingequation}.
The weights are smooth, and the full incoming limit gives
\[
 |w_i(u)-\log h_B(e_i)|
 \le |p_i|\int_{-\infty}^u\sech s\,\dd s
 \le2|p_i|e^u.
\]
Apply Lemma~\ref{lem:taildescent} with time reversed. On a shorter
tail one fixed $K$-basis works at both vertices. The identity
arrow is an isometry there, so equal momentum blocks have
isomorphic reductions and carry precisely the claimed $B$ labels.
\end{proof}

Both end estimates are $O(|t|^{-1})$ in the original coordinate.
The unresolved part of the construction is now confined to a
compact interval. The next section proves that it too has a finite
algebraic description.
